\documentclass[a4paper,12pt]{article}
\usepackage[T2A]{fontenc}
\usepackage[utf8]{inputenc}
\usepackage[russian]{babel}
\usepackage{amsmath,amssymb,mathtools}
\usepackage{helvet}
\renewcommand{\familydefault}{\sfdefault}
\usepackage{icomma}
\usepackage{geometry}
\geometry{left=21mm,right=21mm,top=23mm,bottom=22mm}
\usepackage{microtype}
\usepackage{tikz}
\usetikzlibrary{calc,arrows.meta,positioning,shapes.geometric}
\usepackage{booktabs,tabularx,array}
\usepackage{enumitem}
\usepackage{fancyhdr}
\usepackage{parskip}
\usepackage{setspace}
\usepackage{xcolor}
\usepackage{needspace}
\usepackage{hyperref}
\hypersetup{hidelinks,pdfauthor={Лёвин Артём Александрович},pdftitle={Чек-лист. Классическая вероятность и комбинаторика}}
\definecolor{accent}{HTML}{008080}
\definecolor{darkaccent}{HTML}{004D4D}
\definecolor{textgray}{HTML}{333333}
\definecolor{quiet}{HTML}{637373}
\setlength{\parindent}{0pt}
\setlength{\parskip}{5pt}
\onehalfspacing
\setlength{\headheight}{15pt}
\emergencystretch=3em
\pagestyle{fancy}
\fancyhf{}
\fancyhead[L]{\small\color{quiet}Теория вероятностей и комбинаторика}
\fancyhead[R]{\small\color{quiet}Марк Гукин}
\fancyfoot[C]{\small\color{quiet}\thepage}
\renewcommand{\headrulewidth}{0.25pt}
\renewcommand{\footrulewidth}{0pt}
\newcommand{\Section}[1]{\Needspace{4\baselineskip}\vspace{3pt}{\large\bfseries\color{darkaccent}#1}\par\vspace{2pt}}
\newcommand{\Key}[1]{\textbf{\color{darkaccent}#1}}
\newcommand{\Mark}[1]{\textcolor{accent}{\textbf{#1}}}
\renewcommand{\arraystretch}{1.23}
\setlist[enumerate]{leftmargin=10mm,itemsep=6mm,topsep=3mm}

\begin{document}
\thispagestyle{empty}
\begin{center}
\vspace*{23mm}
{\Large\color{darkaccent} МАТЕМАТИКА\par}
\vspace{16mm}
{\LARGE Чек-лист\par}
\vspace{12mm}
{\LARGE\color{accent} Классическая вероятность\par}
\vspace{1mm}
{\LARGE\color{accent} и комбинаторика\par}
\vspace{16mm}
{\large Подготовка к ЕГЭ по профильной математике\par}
\vfill
{\large Марк Гукин\par}
\vspace{10mm}
{\large 15 ноября 2026 г.\par}
\vspace{15mm}
{Лёвин Артём Александрович\par}
{эксклюзивно для Марка Гукина\par}
\vspace{14mm}
\begin{tikzpicture}[x=1cm,y=1cm]
\draw[accent,line width=.75pt] (0,0) .. controls (2,.8) and (3.2,-.54) .. (5.2,0)
 .. controls (7.3,.77) and (8.2,-.53) .. (10.4,0);
\fill[darkaccent] (5.2,0) circle (1.05pt);
\end{tikzpicture}
\end{center}
\clearpage

\Section{1. Вероятность: благоприятные и все исходы}
\Key{Классическая формула вероятности} применяется, когда все элементарные исходы случайного опыта равновозможны:
\[
\boxed{P(A)=\frac{M}{N}},\qquad 0\le M\le N,\quad N>0.
\]
Здесь $N$ --- число всех элементарных исходов, $M$ --- число исходов, при которых происходит событие $A$.

\textbf{Пример 1.} Случайно выбирают одну цифру от $0$ до $9$. Какова вероятность получить чётную цифру?
\[
\underbrace{\{0,2,4,6,8\}}_{M=5},\qquad N=10,
\qquad P=\frac5{10}=\boxed{\frac12}.
\]
Нуль --- чётное число: он делится на $2$ без остатка.

\Key{Противоположное событие:} если событие $\overline A$ означает «$A$ не произошло», то
\[
P(A)=1-P(\overline A),\qquad M_A=N-M_{\overline A}.
\]
Такой подсчёт удобен, когда исключений меньше, чем благоприятных случаев.

\Section{2. Два контрольных вопроса комбинаторики}
Перед подсчётом вариантов всегда определяйте:
\begin{enumerate}[label=\textbf{\arabic*.},itemsep=1mm,topsep=1mm]
\item \Key{Допускаются ли повторения?} Может ли один объект быть выбран несколько раз?
\item \Key{Важен ли порядок?} Различаются ли результаты после перестановки выбранных объектов?
\end{enumerate}

\begin{center}
\begin{tabularx}{\linewidth}{@{}p{.29\linewidth}XX@{}}
\toprule
Ситуация & Порядок & Повторения\\
\midrule
Два человека занимают первое и второе места & Важен & Запрещены\\
Двое образуют одну команду & Не важен & Запрещены\\
Два человека независимо выбирают цифру & Важен: кто что выбрал & Допускаются\\
\bottomrule
\end{tabularx}
\end{center}

\Key{Метод позиций.} Если последовательно есть $n_1,n_2,\ldots,n_k$ допустимых вариантов, число исходов равно $n_1n_2\cdots n_k$.
Например, в очереди из трёх различных человек: $3\cdot2\cdot1=3!=6$ исходов.

\Section{3. Когда необходима формула сочетаний}
Если из $n$ различных объектов выбирают $k$ объектов \emph{без повторений и без учёта порядка}, то
\[
\boxed{\binom nk=\frac{n!}{k!(n-k)!}},\qquad 0\le k\le n.
\]
Знак $!$ означает факториал: $4!=4\cdot3\cdot2\cdot1$, $0!=1$. Сначала вычисляйте разность $n-k$, затем факториал $(n-k)!$.

\Section{4. Выбор группы и фиксированный участник}
\textbf{Пример 2.} Из пяти туристов наугад выбирают двоих для похода в магазин. Какова вероятность, что выбран турист А?

Всего неупорядоченных пар $N=\binom52=10$. При условии, что А уже выбран, второго можно выбрать среди четырёх остальных: $M=\binom41=4$.
\[
P=\frac{\binom41}{\binom52}=\frac4{10}=\boxed{\frac25}.
\]
\Key{Важно:} пару $\{\text{А},\text{Б}\}$ и пару $\{\text{Б},\text{А}\}$ учитывают один раз.

\Section{5. Самопроверка таблицей}
Два человека А и Б \emph{независимо} выбирают одну цифру из $\{1,2,3\}$. Результат «А выбрал $1$, Б выбрал $2$» отличается от обратного результата.
\begin{center}
\begin{tabular}{@{}c ccc@{}}
\toprule
& \multicolumn{3}{c}{Цифра человека Б}\\
Цифра А & $1$ & $2$ & $3$\\
\midrule
$1$ & $(1,1)$ & $(1,2)$ & $(1,3)$\\
$2$ & $(2,1)$ & $(2,2)$ & $(2,3)$\\
$3$ & $(3,1)$ & $(3,2)$ & $(3,3)$\\
\bottomrule
\end{tabular}
\end{center}
Всего $N=3\cdot3=9$ равновозможных исходов. Совпадений ровно три, поэтому
\[
P(\text{разные цифры})=\frac{9-3}{9}=\boxed{\frac23}.
\]
Проверка методом позиций: $M=3\cdot2=6$. При независимом выборе порядок участников \emph{существенен}; деление на $2!$ исказит подсчёт.

\Needspace{12\baselineskip}
\Section{6. Условие «рядом» и метод объединения}
\textbf{Пример 3.} Четыре различных человека случайно садятся в ряд. Какова вероятность, что А и Б сидят рядом?

Все рассадки: $N=4!=24$. Считайте пару А--Б одним блоком; вместе с двумя другими людьми получится $3!$ способов. Внутри блока люди меняются местами двумя способами:
\[
M=3!\cdot2!=12,\qquad P=\frac{12}{24}=\boxed{\frac12}.
\]
\Key{Не забудьте:} если важен порядок внутри блока, внутренние перестановки обязательно учитываются.

\Section{7. Разбиение на пары}
Если четырёх различных участников делят на две \emph{неразличимые} пары, достаточно выбрать напарника для А. Возможны три выбора, после чего оставшаяся пара определяется автоматически. Вероятность, что А будет с Б, равна $\boxed{1/3}$.

\Section{8. Вероятность выбрать диагональ многоугольника}
Две различные вершины $n$-угольника выбирают случайным образом, без учёта порядка. Всего соединяющих их отрезков $\binom n2$. Среди них $n$ сторон, все остальные --- диагонали:
\[
D=\binom n2-n=\frac{n(n-3)}2.
\]
Значит, для $n\ge3$ вероятность получить диагональ равна
\[
P=\frac{\binom n2-n}{\binom n2}=\boxed{\frac{n-3}{n-1}}.
\]
\textbf{Пример 4.} Для квадрата $P=(4-3)/(4-1)=\boxed{1/3}$. Вычитать число сторон полезнее, чем перечислять диагонали большого многоугольника.

\Section{9. Независимые испытания: ровно $k$ успехов}
В каждом из $n$ независимых испытаний возможны два равновероятных результата: «успех» и «неудача». Тогда число всех последовательностей $2^n$, а число способов выбрать $k$ позиций успеха равно $\binom nk$:
\[
\boxed{P(\text{ровно }k\text{ успехов})=\frac{\binom nk}{2^n}}.
\]
\Needspace{9\baselineskip}\textbf{Пример 5.} В трёх независимых подбрасываниях честной монеты нужно получить ровно два орла. Подходят ООР, ОРО, РОО. Поэтому
\[
P=\frac{\binom32}{2^3}=\boxed{\frac38}.
\]
Выбираются \emph{позиции} успехов. Если вероятность успеха отличается от $1/2$, эту формулу в таком виде применять нельзя.

\Section{10. Кратность четырём: метод двух последних цифр}
Натуральное число делится на $4$ тогда и только тогда, когда число из его двух последних цифр делится на $4$.

\textbf{Пример 6.} Из цифр $1,2,3,4,5,6$ составляют шестизначное число, каждую цифру используя ровно один раз. Сколько чисел кратно $4$?

Подходящие окончания из двух \emph{разных} цифр:
\[
12,\quad16,\quad24,\quad32,\quad36,\quad52,\quad56,\quad64.
\]
Окончаний восемь. Оставшиеся четыре цифры можно переставить $4!$ способами:
\[
\boxed{8\cdot4!=192\text{ числа}}.
\]
Если среди допустимых цифр есть нуль, проверяйте ведущую цифру: натуральное шестизначное число не может начинаться с $0$.

\Section{11. Типичные ошибки}
\begin{itemize}[leftmargin=5mm,itemsep=1mm,topsep=1mm]
\item Делить на $k!$ по привычке, хотя перестановка объектов меняет исход.
\item Учитывать пару $\{\text{А},\text{Б}\}$ дважды при выборе неупорядоченной группы.
\item Путать количество \emph{исходов} и вероятность: $M$ и $N$ --- целые числа, $P$ может быть дробью.
\item Применять $M/N$ без проверки равновозможности элементарных исходов.
\item Забывать внутренние перестановки объединённого блока и ограничение на начальный нуль.
\end{itemize}
\clearpage

\begin{center}
{\Large\bfseries\color{darkaccent}Алгоритм: как считать классическую вероятность}
\end{center}
\vspace{4mm}
\begin{center}
\begin{tikzpicture}[
  >={Latex[length=2.5mm]}, line width=.75pt, font=\small,
  start/.style={ellipse,draw=accent,text width=41mm,align=center,minimum height=11mm,inner sep=2pt},
  action/.style={rectangle,rounded corners=2pt,draw=accent,text width=58mm,align=center,minimum height=13mm,inner sep=3pt},
  branch/.style={rectangle,rounded corners=2pt,draw=accent,text width=52mm,align=center,minimum height=14mm,inner sep=3pt},
  decision/.style={diamond,aspect=2.5,draw=accent,text width=32mm,align=center,inner sep=1.8pt},
  link/.style={->,draw=darkaccent}
]
\node[start] (s) at (0,0) {Начало};
\node[action] (model) at (0,-1.8) {Выделите равновозможные элементарные исходы};
\node[decision] (order) at (0,-4.05) {Важен порядок?};
\node[branch] (ordered) at (-3.45,-6.45) {Да: учитывайте разные последовательности};
\node[branch] (unordered) at (3.45,-6.45) {Нет: считайте группы без перестановок};
\node[action] (repeats) at (0,-8.8) {Установите, допустимы ли повторения, и учтите ограничения};
\node[action] (count) at (0,-11.25) {Найдите число всех исходов $N$ и благоприятных $M$};
\node[decision] (check) at (0,-13.75) {Верно ли $0\le M\le N$ и $N>0$?};
\node[branch] (repair) at (4.9,-15.95) {Нет: пересмотрите подсчёт и модель исходов};
\node[start] (end) at (0,-17.25) {$P(A)=M/N$};
\draw[link] (s) -- (model);
\draw[link] (model) -- (order);
\draw[link] (order.west) -| node[above,pos=.27,fill=white,inner sep=1pt]{да} (ordered.north);
\draw[link] (order.east) -| node[above,pos=.27,fill=white,inner sep=1pt]{нет} (unordered.north);
\draw[link] (ordered.south) |- (repeats.west);
\draw[link] (unordered.south) |- (repeats.east);
\draw[link] (repeats) -- (count);
\draw[link] (count) -- (check);
\draw[link] (check.south) -- node[left,fill=white,inner sep=1pt]{да} (end.north);
\draw[link] (check.east) -| node[above,pos=.25,fill=white,inner sep=1pt]{нет} (repair.north);
\draw[link] (repair.east) -- ++(.25,0) |- ($(count.east)+(2.6,0)$) -- (count.east);
\end{tikzpicture}
\end{center}
\clearpage

{\Large\bfseries\color{darkaccent}Тренировочный блок}\par
\vspace{2mm}
\textcolor{quiet}{Путь А · 10 заданий · от базового уровня к комплексному. В задачах со случайным выбором допустимые исходы считаются равновозможными, если не сказано иное.}
\vspace{2mm}
\begin{enumerate}[label=\textbf{\arabic*.}]
\item Случайно выбирают одну цифру от $0$ до $9$. Найдите вероятность того, что выбранная цифра кратна $3$.

\item Четыре различных человека случайно становятся в очередь. Какова вероятность, что Анна стоит раньше Бориса?

\item Из семи добровольцев равновероятно выбирают группу из трёх человек. Какова вероятность, что в группу попадёт Марк?

\item Из цифр $1,2,3,4,5$ составляют пятизначные числа, используя каждую цифру ровно один раз. Сколько таких чисел делится на $4$?

\item Из двенадцати вершин правильного двенадцатиугольника равновероятно выбирают две различные вершины. Какова вероятность, что соединяющий их отрезок --- диагональ?

\item Пять разных человек случайно рассаживаются в один ряд из пяти стульев. Какова вероятность, что два заранее названных человека сидят рядом?

\item Два участника А и Б независимо друг от друга выбирают по одной цифре из $\{1,2,3,4\}$. Какова вероятность, что выбранные цифры различны?
\end{enumerate}
\clearpage
{\Large\bfseries\color{darkaccent}Тренировочный блок · продолжение}\par
\vspace{4mm}
\begin{enumerate}[label=\textbf{\arabic*.},start=8]
\item Шестерых разных участников случайно делят на три пары. Номера пар не имеют значения. Какова вероятность, что Анна и Борис окажутся в одной паре?

\item Честную монету независимо подбрасывают пять раз. Найдите вероятность того, что орёл выпадет ровно три раза.

\item Из цифр $0,1,2,3,4,5$ случайно составляют шестизначное число, используя каждую цифру ровно один раз; начальный нуль запрещён. Найдите вероятность того, что полученное число делится на $4$.
\end{enumerate}
\clearpage

{\Large\bfseries\color{darkaccent}Ответы}\par
\vspace{7mm}
\begin{center}
\begin{tabular}{@{}p{.16\linewidth}p{.66\linewidth}@{}}
\toprule
№ & Ответ\\
\midrule
1 & $\dfrac25$\\
2 & $\dfrac12$\\
3 & $\dfrac37$\\
4 & $24$ числа\\
5 & $\dfrac9{11}$\\
6 & $\dfrac25$\\
7 & $\dfrac34$\\
8 & $\dfrac15$\\
9 & $\dfrac5{16}$\\
10 & $\dfrac6{25}$\\
\bottomrule
\end{tabular}
\end{center}
\vspace{10mm}
\textcolor{quiet}{Проверьте сначала число всех исходов, затем число благоприятных. Для вероятности ответ должен находиться в промежутке от $0$ до $1$.}
\end{document}
